\section{Quy trình thực hiện và tối ưu}

Mã mẫu trong thư mục \texttt{legacy} được giữ nguyên để đối chiếu. Phần
triển khai mới được tổ chức thành các mô-đun \texttt{Board}, \texttt{Piece},
\texttt{Game}, \texttt{Input}, \texttt{Renderer} và \texttt{Score} trong
thư mục \texttt{src}. Cách tổ chức này tách kiểm tra va chạm, điều khiển và
hiển thị, trong khi \texttt{Game} điều phối trạng thái của khối đang rơi.

Vòng lặp game kiểm tra phím khoảng mỗi 16 ms, còn tốc độ rơi tự động được
điều khiển bởi bộ đếm thời gian riêng với chu kỳ 500 ms. Vì vậy, các phím
\texttt{A}, \texttt{D}, \texttt{W} và \texttt{X} được phản hồi độc lập với
tốc độ rơi. Bộ kết xuất chỉ tạo khung hình khi trạng thái thay đổi và đưa
con trỏ về đầu màn hình thay vì gọi \texttt{system("cls")} ở mỗi vòng lặp.
Điều này giảm nhấp nháy, giảm chi phí gọi tiến trình hệ thống và tránh độ
trễ do chờ chu kỳ rơi.

\subsection{Kiểm thử}

Sau khi cài MSYS2 UCRT64, mã nguồn được biên dịch bằng:

\begin{verbatim}
g++ -std=c++17 -Wall -Wextra src/*.cpp -o build/tetris.exe
\end{verbatim}

Các kịch bản kiểm thử gồm di chuyển đến sát tường, di chuyển vào khối đã
khóa, rơi nhanh, xoay cạnh tường, từ chối xoay khi mọi vị trí wall-kick đều
không hợp lệ, khóa khối, xóa dòng, tính điểm và thoát bằng phím \texttt{Q}.
Tệp thực thi trong thư mục \texttt{build} là sản phẩm biên dịch cục bộ và
được loại khỏi Git bằng quy tắc \texttt{build/*.exe}.
